%%----------------------------------------------------------
%% proposal_temp.tex
%% v1.0
%% 2023/1/15
%% by XIE Yuhao, e-mail: yuhaox@akane.waseda.jp
%%----------------------------------------------------------

\documentclass{wsdpp}

\title{Waseda Project Proposal Template}
\author{XXX, XXXXX (Student ID: XXXXXXXX-X)}
\affil{\textit{XXXXXX Laboratory}\\
\textit{Graduate School of Information, Production and Systems, Waseda University}\\
\textit{Fukuoka, Japan}\\
E-mail: XXXXX@XXXXX.waseda.jp}
\date{15\textsuperscript{th} January, 2023}

\begin{document}

\maketitle

\tableofcontents

\setcounter{page}{0}

\chapter{Chapter}

\lettrine{D}{ropped} caps at the start of the first paragraph. Make sure the text you put in the first paragraph exceeds two lines, otherwise the dropped caps will not look as good as intended.

Indent begins on second paragraph.
Dummy text: \lipsum[1]

\section{Section}
Dummy text: \lipsum[1]

\subsection{subsection}
Dummy text: \lipsum[1]

\subsubsection{subsubsection}
Dummy text: \lipsum[1]

\chapter{Floats}
\section{Figures and Minipages}
An example of a figure is shown in Figure \ref{fig:example}. The recommended width of the figure is $0.9\times$\textbackslash linewidth, and the recommended width:height aspect ratio is $1.41:1$ (i.e. A4 paper)

Additionally, one can create a minipage for extra information like biography as shown in the following:

\footnotetext[1]{acquired from the official website of IPS at\\https://www.waseda.jp/fsci/gips/en/about/overview-2/}

Dummy text: \lipsum[1]

\begin{figure}[b!]
    \centering
    \includegraphics[width=.5\linewidth]{logo_gips_unit.png}
    \caption{The logo of IPS, Waseda University.}
    \label{fig:example}
\end{figure}

\begin{figure}[b!]
\fcolorbox{wsdred}{wsdgrey}{ 
%% first argument is the frame color
%% second argument is the background color
\begin{minipage}{.95\linewidth}
\color{white}

%% the title of the minipage
\begin{center}
    \fontsize{10}{12}\fontfamily{phv}\fontshape{sc}\selectfont
    \textbf{About the School}
\end{center}

\begin{wrapfigure}{r}{.3\linewidth}
\includegraphics[width=\linewidth]{logo_gips_unit.png}
\end{wrapfigure}
\textbf{Waseda University Graduate School of Information, Production and Systems (IPS)} is a graduate school that has no corresponding undergraduate department established within the Kitakyushu Science and Research Park area in 2003 as the base for Waseda University to expand its presence in Asia. With three fields of study, including Information Architecture, Production System, and Integrated Systems, IPS undertakes academic research in the technological fields that society currently requires, and strives to attain a sustainable society through the use of technology. \footnotemark[1].
\end{minipage}
}
\end{figure}

\section{Tables}
Creating tables is the same as any other latex documents, as shown in Table \ref{tab:example}.

Dummy text: \lipsum[1]

\begin{table}[b!]
    \centering
    \begin{tabular}{ccc}
    \hline
        header1 &  header2 & header 3\\
        \hline
        data1 & data2 & data3\\
        data4 & data5 & data6
    \end{tabular}
    \caption{An example of a floating table.}
    \label{tab:example}
\end{table}

\chapter{Equations}
Just normally input your equations as either inline one: $E=h\frac{\lambda}{c}$ or with equation environment (equation* to remove numbering):

\begin{equation}
\left\{
\begin{aligned}
    &\oiint_{\partial\Omega}{\mathbf{E}\cdot d\mathbf{S}}=\frac 1{\varepsilon_0}\iiint_\Omega{\rho dV}\\
    &\oiint_{\partial\Omega}{\mathbf{B}\cdot d\mathbf{S}}=0\\
    &\oint_{\partial\Sigma}{\mathbf{E}\cdot d\boldsymbol{l}}=-\frac d{dt} \iint_\Sigma{{\mathbf{B}\cdot d\mathbf{S}}}\\
    &\oint_{\partial\Sigma}{\mathbf{B}\cdot d\boldsymbol{l}}=\mu_0 \left(\iint_\Sigma{{\mathbf{J}\cdot d\mathbf{S}}}+\varepsilon_0\frac d{dt} \iint_\Sigma{{\mathbf{E}\cdot d\mathbf{S}}}\right)\cite{Maxwell1865}
\end{aligned}
\right.
\end{equation}

\newpage % added here to prevent an issue where \appendices causes the right header to appear as "Appendix" at the last page of normal chapters

\begin{appendices}
\chapter{}
Appendix one text goes here.

% you can choose not to have a title for an appendix
% if you want by leaving the argument blank
\chapter{}
Appendix two text goes here.

\newpage % same as before, but the other way around
\end{appendices}

% use chapter* for acknowledgment
\chapter*{Acknowledgment}
\rhead{\small\sffamily\bfseries Acknowledgment}
The authors would like to thank...

\newpage



\bibliographystyle{IEEEtran}
%% you can change the style into any other styles available, I personally love IEEEtran.
\bibliography{citations}
%% to generate references, input the name of your .bib file and cite anywhere in the document.
\rhead{\small\sffamily\bfseries \bibname}

\end{document}
